\chapter{无源通信：散射变换与波变量}\label{ch:teleop-wave}

% ════════════════════════════════════════════════════════════════
%  机械臂部·遥操作子部 第三章。承 ch:teleop-twoport（二端口/透明度）、
%  ch:teleop-stability（Llewellyn/时延使 η<0）；本章给 Anderson–Spong(1989)
%  散射变换 + Niemeyer–Slotine(1991) 波变量：用信号变换让通道在任意常数时
%  延下天然无源。引出 ch:teleop-tdpa（时域无源/能量罐）。记号归一到本书：
%  右扰动主线、tau=Jᵀf、端口功率 P=fᵀv（与 \cref{eq:os-port-power} 一致）。
% ════════════════════════════════════════════════════════════════

\begin{practice}[前置自测]
开始之前，先试答下列问题。若已能流畅作答，可只速读本章的散射矩阵与波变量速查卡；若卡住，正是本章要为你解决的。
\begin{enumerate}
  \item \textbf{无源性的能量表述}：写出单端口系统无源的能量不等式。其中初始储能项代表什么？（提示：\cref{eq:os-passive-integral}）
  \item \textbf{时延的破坏}：上一章给出理想透明时 Llewellyn 稳定因子 $\eta=0$；通信延迟 $T$ 如何使 $\eta<0$？关键频率 $\omega=\pi/T$ 处发生了什么？
  \item \textbf{传输线}：无损 $LC$ 传输线的波速 $c=1/\sqrt{LC}$ 与特征阻抗 $Z_0=\sqrt{L/C}$ 各是什么含义？
  \item \textbf{d'Alembert 解}：一维波动方程通解里的正向波与反向波各代表什么物理量的传播？
  \item 直接把主端力\emph{延迟}后送给从端（$f_s(t)=f_m(t-T)$），为什么会“凭空产生能量”？
\end{enumerate}
\end{practice}

\section*{本章目标}
学完本章，你应当能够：
\begin{enumerate}
  \item \textbf{说清}为何“延迟保能、放大产能”，从而把遥操作稳定性问题从\emph{控制问题}重述为\emph{信号表示问题}；
  \item \textbf{从传输线方程推导}散射变换，并理解散射矩阵的 $\|\mathbf S\|_\infty\le1$ 与无源性的等价；
  \item \textbf{掌握}波变量的完整定义、逆变换与功率恒等式 $f\dot x=\tfrac12(u^2-v^2)$，并把它写成功率保持变换 $\boldsymbol\Phi^\top\boldsymbol\Sigma\boldsymbol\Phi=\mathbf J$；
  \item \textbf{推导} Niemeyer--Slotine 主/从控制律 $f_s=\sqrt{2b}\,u_s-b\,\dot x_s$，并用三段分解复用互联定理（\cref{thm:os-interconnect}）证明全系统无源；
  \item \textbf{独立完成}任意常数时延下通道无源性的四步证明，并解释为何时延 $T$ \emph{不}进入非负条件；
  \item \textbf{选择}波阻抗 $b$：理解自由空间残余\emph{惯量} $M_{\rm eff}=bT$（注意传递阻抗 $Z_{\rm to}=\mathrm j\omega\,bT$ 是\emph{质量}型而非阻尼型）、刚墙渲染刚度 $K_{\rm eff}=b/T$、阻抗匹配 $b=\sqrt{K_eM_m}$ 与振铃；
  \item \textbf{识别}波变量的位置漂移问题，并知道它由波积分（wave integral）修复（详见 \cref{ch:teleop-tdpa}）。
\end{enumerate}

\subsection*{本章知识导航}
本章是一条“\emph{把力/速度信号搬进波域，让通信通道天然无源}”的主线。结构如下：

\begin{center}
\small
\begin{tikzpicture}[
  node distance=4.5mm and 7mm, >=Stealth,
  blk/.style={draw, rounded corners, align=center, font=\small, inner sep=3pt, fill=main!6, draw=main!60!black},
  grp/.style={draw, rounded corners, align=center, font=\small, inner sep=3pt, fill=second!6, draw=second!60!black},
  ln/.style={->, thick, gray!60!black}]
\node[grp] (mot) {动机\\ 时延使 $\eta<0$};
\node[blk, right=of mot] (tl) {传输线\\ 行波/功率分解};
\node[blk, right=of tl] (sc) {散射变换\\ $\boldsymbol\Phi,\;\|\mathbf S\|_\infty\!\le\!1$};
\node[blk, below=7mm of mot] (wv) {波变量\\ Niemeyer--Slotine};
\node[blk, right=of wv] (pf) {时延无源\\ 四步证明};
\node[grp, right=of pf] (b) {波阻抗 $b$\\ 匹配/振铃};
\node[grp, below=7mm of wv, fill=third!8, draw=third!60!black] (drift) {位置漂移\\ wave integral};
\draw[ln] (mot)--(tl); \draw[ln] (tl)--(sc);
\draw[ln] (sc.south) -- ++(0,-3.6mm) -| (wv.north);
\draw[ln] (wv)--(pf); \draw[ln] (pf)--(b);
\draw[ln] (wv.south)--(drift.north);
\end{tikzpicture}
\end{center}

\noindent\textbf{推荐路径}：第 1--3 节建立物理直觉与散射变换的代数地基（任何读者的主干）；第 4--5 节是本章\emph{核心}——波变量控制律与任意常数时延的无源性证明；第 6 节把抽象的 $b$ 落到“多软的墙、多黏的自由空间”这类可感量上；第 7 节指出残留的工程问题并把它转交下一章。

\subsection*{前置知识桥接}
\cref{ch:teleop-twoport} 已把主从遥操作抽象成\emph{二端口网络}，给出混合参数 $\mathbf H$ 与透明度的定义；\cref{ch:teleop-stability} 用 Llewellyn 准则揭示了“透明度--稳定性”的根本矛盾，并证明通信延迟会把稳定因子 $\eta$ 压到负值。本章要回答的是上一章\emph{留下}的问题：能否不动控制器、只换\emph{传输的信号形式}，让通道在任意常数时延下天然无源？所需的无源性语言——端口功率 $P=\mathbf f^\top\mathbf v$（\cref{eq:os-port-power}）、耗散不等式（\cref{eq:os-passive-integral}）、被动系统互联定理（\cref{thm:os-interconnect}）、能量罐（\cref{thm:os-energytank}）——均已在操作空间章 \cref{ch:operational-space} 严格建立，本章直接复用、不再重证。

\subsection*{如果跳过本章会怎样}
\begin{itemize}
  \item 你将无法理解 \cref{ch:teleop-tdpa} 的时域无源控制（TDPA）——它本质是波变量思想在\emph{时域、对时变时延}的推广，不懂波变量就只能照抄能量罐代码而不会调试；
  \item 面对“100\,ms 往返时延下主从一接触硬墙就发散”的真实故障，你将只会盲目加阻尼（牺牲透明度），而不知道存在一种\emph{结构性}的解法。
\end{itemize}

\begin{note}
\textbf{本章记号与约定.}\;
力/力矩用 $f$、广义力对偶仍守 $\boldsymbol\tau=\mathbf J^\top\mathbf f$；速度一律写 $\dot x$（标量自由度）或 $\mathbf v$，端口功率 $P=f\dot x$ 与 \cref{eq:os-port-power} 一致。
\textbf{端口方向}：本章端口功率以“流入通信通道”为正。主端取流入速度 $\nu_m=\dot x_m$；从端物理速度 $\dot x_s$ 通常由主端指令流向环境，故对通道而言流入速度 $\nu_s=-\dot x_s$。
\textbf{波变量}用 $u$（正向波，主$\to$从）、$v$（反向波，从$\to$主），波阻抗 $b>0$；为避免与端口\emph{速度}混淆，速度从不写作 $v$。
源文献 \cite{anderson1989bilateral,niemeyer1991stable} 常用 $[F_h;-V_s]=\mathbf H[V_m;F_e]$ 的大写端口变量与 $f\!_v=(u^2-v^2)/2$ 记法、并以 $\lambda_p$ 标度位置 $x_s=\lambda_p x_m$；本章统一转换到上述小写记号，差异处以批注说明。
\end{note}

% ════════════════════════════════════════════════════════════════
\section{动机：延迟为何能让“无源”失守}
\label{sec:tw-motivation}

\paragraph{反面教训：直接传力会凭空造能量。}
最朴素的双边遥操作是“力--力”或“位置--力”直传：从端跟随主端运动，主端反馈从端测得的环境力。一旦通信有延迟 $T$，从端施加的是\emph{过时}的主端力。考虑从端瞬时功率
\begin{equation}
  P_s(t)=f_m(t-T)\,\dot x_s(t).
  \label{eq:tw-naive-power}
\end{equation}
$f_m(t-T)$ 与 $\dot x_s(t)$ 是\emph{两个不同时刻}的信号，其乘积可正可负。当力与速度因延迟而“失相”——例如 $T$ 时刻前的力为正、当前速度也为正——乘积为正，等效于通道\emph{向系统注入}能量。上一章 \cref{ch:teleop-stability} 已从频域给出同一结论：延迟项 $\mathrm e^{-\mathrm j\omega T}$ 使 $h_{12}h_{21}$ 的实部在 $\omega=\pi/T$ 附近变号，Llewellyn 稳定因子 $\eta(\omega)<0$，理想透明的通道由\emph{无源}沦为\emph{有源}，接触硬物时主从回路自激发散。

\paragraph{核心洞察（Anderson--Spong 1989）：延迟保能，放大才产能。}
解决的钥匙来自一个物理事实\cite{anderson1989bilateral}：\textbf{延迟本身不产生能量}。电磁波沿无损同轴电缆传播 $100\,$m，到达时间推迟约 $0.5\,\mu$s，但幅值分毫未改——能量只是“晚到”，并未增减。能量守恒由物理定律保证。真正产能的是\emph{放大}：把信号幅值放大需要外部电源，否则违反热力学第一定律。于是问题转化为：

\begin{insight}
若能找到一种信号表示，使通信通道的作用是\emph{纯延迟}而非放大，则通道自动无源——无论时延多大。\textbf{遥操作的稳定性问题，被重述为一个信号表示问题}：不是改控制器，而是改传输的信号形式。
\end{insight}

\paragraph{不是让延迟消失，而是让延迟变安全。}
须先破一个误解：波变量\emph{不}消除延迟。操作者仍会感到延迟带来的透明度下降——延迟的波到达后解码出的力/速度终归是“过时的”。波变量做的是把延迟\emph{从“不安全”变为“安全”}：稳定性与透明度是两个独立指标，本章解决前者，后者仍受延迟约束（\cref{sec:tw-impedance} 量化）。这一“先把系统救活、再谈手感”的分工，是后续所有时延遥操作方法的共同起点。

% ════════════════════════════════════════════════════════════════
\section{传输线：从行波到功率分解}
\label{sec:tw-transmissionline}

要把通道做成“纯延迟”，先看自然界里现成的纯延迟器件——无损传输线。

\paragraph{无损 $LC$ 线的波动方程。}
设单位长度电感 $L$、电容 $C$，电压 $V(x,t)$、电流 $I(x,t)$ 满足电报方程
\begin{equation}
  \frac{\partial V}{\partial x}=-L\frac{\partial I}{\partial t},\qquad
  \frac{\partial I}{\partial x}=-C\frac{\partial V}{\partial t}.
  \label{eq:tw-telegraph}
\end{equation}
消去 $I$ 得波动方程 $\partial^2_xV=LC\,\partial^2_tV$，波速 $c=1/\sqrt{LC}$，特征阻抗 $Z_0=\sqrt{L/C}$。

\paragraph{d'Alembert 解：正/反两支行波。}
通解为左右两支独立行波之和\cite{anderson1989bilateral}：
\begin{align}
  V(x,t)&=V^+(t-x/c)+V^-(t+x/c),\label{eq:tw-dalembert-v}\\
  I(x,t)&=\tfrac{1}{Z_0}\big[V^+(t-x/c)-V^-(t+x/c)\big],\label{eq:tw-dalembert-i}
\end{align}
其中 $V^+$ 是从左向右传播的\emph{正向}行波，$V^-$ 是从右向左的\emph{反向}行波。

\begin{derivation}[功率分解：传输线无源的根因]
把 \cref{eq:tw-dalembert-v,eq:tw-dalembert-i} 代入瞬时功率 $P=VI$，省略参数后
\begin{equation}
  P=VI=\frac{1}{Z_0}\big(V^++V^-\big)\big(V^+-V^-\big)
   =\underbrace{\frac{(V^+)^2}{Z_0}}_{P^+\ \text{正向波功率}}-\underbrace{\frac{(V^-)^2}{Z_0}}_{P^-\ \text{反向波功率}}.
  \label{eq:tw-tl-power}
\end{equation}
\end{derivation}

\cref{eq:tw-tl-power} 是本章一切的种子：传输线功率是\emph{两个独立方向波功率之差}。两支波独立传播、互不干扰——正向波被延迟不影响反向波。能量被延迟传递，但既不被创造也不被消灭，这正是无损传输线恒无源的原因。

\begin{table}[H]\centering\small
\caption{电气--力学--波变量三重类比}
\label{tab:tw-analogy}
\begin{tabular}{lll}
\toprule
电气传输线 & 力学遥操作 & 波变量 \\
\midrule
电压 $V$ & 力 $f$ & --- \\
电流 $I$ & 速度 $\dot x$ & --- \\
正向波 $V^+$ & --- & 正向波 $u$ \\
反向波 $V^-$ & --- & 反向波 $v$ \\
特征阻抗 $Z_0$ & --- & 波阻抗 $b$ \\
功率 $\dfrac{(V^+)^2-(V^-)^2}{Z_0}$ & $f\dot x$ & $\dfrac12\big(u^2-v^2\big)$ \\
\bottomrule
\end{tabular}
\end{table}

\begin{insight}
散射变换与傅里叶变换在精神上同源：傅里叶把卷积变乘法，散射把“力$\times$速度”的功率变成“两个独立非负项之差”。二者都是用一次坐标变换，让某个守恒律在新坐标下\emph{显式可见}。这一历史脉络也耐人寻味——散射参数本是 $20$ 世纪 $40$ 年代微波工程师描述天线/滤波器端口的工具，因遥操作通道与微波传输线\emph{数学同构}（都是两端口、都涉双向传播与延迟、都须保证不放大），Anderson 与 Spong 把它整体迁移了过来。
\end{insight}

% ════════════════════════════════════════════════════════════════
\section{散射变换：把功率对角化}
\label{sec:tw-scattering}

\paragraph{定义。}
引入波阻抗 $b>0$（类比 $Z_0$），定义入射波 $a$（流入网络的能量载体）与反射波 $r$（流出）：
\begin{equation}
  a=\frac{f+b\dot x}{\sqrt{2b}},\qquad
  r=\frac{f-b\dot x}{\sqrt{2b}}.
  \label{eq:tw-scatter-def}
\end{equation}
归一化因子 $1/\sqrt{2b}$ 经精心选取，使 $a^2,r^2$ 直接具有功率量纲：$[f]=$N、$[b\dot x]=$(Ns/m)(m/s)$=$N、$[\sqrt{2b}]=\sqrt{\text{Ns/m}}$，故 $[a^2]=\text{N}^2/(\text{Ns/m})=\text{N}\!\cdot\!\text{m/s}=$W。

\paragraph{逆变换。}
\cref{eq:tw-scatter-def} 是 $2\times2$ 线性映射，可逆：
\begin{equation}
  f=\sqrt{\tfrac{b}{2}}\,(a+r),\qquad
  \dot x=\frac{a-r}{\sqrt{2b}}.
  \label{eq:tw-scatter-inv}
\end{equation}
代回即验 $f\mapsto f$、$\dot x\mapsto\dot x$，无任何信息损失。

\begin{derivation}[功率恒等式]
直接把 \cref{eq:tw-scatter-inv} 相乘：
\begin{equation}
  f\dot x=\sqrt{\tfrac{b}{2}}(a+r)\cdot\frac{a-r}{\sqrt{2b}}
        =\tfrac12(a+r)(a-r)=\boxed{\tfrac12\big(a^2-r^2\big)}.
  \label{eq:tw-power-identity}
\end{equation}
\end{derivation}

\cref{eq:tw-power-identity} 把无源性判据彻底简化。单端口无源要求 $\int_0^t f\dot x\,\mathrm d\tau\ge0$（\cref{eq:os-passive-integral} 的标量形），在波域等价于
\begin{equation}
  \int_0^t a^2(\tau)\,\mathrm d\tau\ \ge\ \int_0^t r^2(\tau)\,\mathrm d\tau,
  \label{eq:tw-passive-wave}
\end{equation}
即入射波总能量不少于反射波总能量——网络\emph{吸收}而非\emph{制造}能量。分析从“力$\times$速度的符号纠缠”降为“两个非负量比大小”。

\paragraph{散射矩阵与功率保持。}
令 $\mathbf p=[f,\dot x]^\top$、$\mathbf w=[a,r]^\top$，\cref{eq:tw-scatter-def} 写成 $\mathbf w=\boldsymbol\Phi\mathbf p$，
\begin{equation}
  \boldsymbol\Phi=\frac{1}{\sqrt{2b}}\begin{bmatrix}1&b\\[1pt]1&-b\end{bmatrix},\qquad
  \det\boldsymbol\Phi=\frac{-2b}{2b}=-1.
  \label{eq:tw-Phi}
\end{equation}
端口功率与波功率分别是两个对称双线性型 $f\dot x=\mathbf p^\top\mathbf J\mathbf p$、$\tfrac12(a^2-r^2)=\mathbf w^\top\boldsymbol\Sigma\mathbf w$，其中
\begin{equation}
  \mathbf J=\tfrac12\begin{bmatrix}0&1\\1&0\end{bmatrix},\qquad
  \boldsymbol\Sigma=\tfrac12\begin{bmatrix}1&0\\0&-1\end{bmatrix}.
  \label{eq:tw-JSigma}
\end{equation}
功率恒等式 $\mathbf p^\top\mathbf J\mathbf p=\mathbf w^\top\boldsymbol\Sigma\mathbf w=\mathbf p^\top\boldsymbol\Phi^\top\boldsymbol\Sigma\boldsymbol\Phi\mathbf p$ 对一切 $\mathbf p$ 成立，故

\begin{proposition}[散射变换是功率保持变换]\label{prop:tw-power-preserving}
\cref{eq:tw-Phi} 的散射矩阵满足
\begin{equation}
  \boldsymbol\Phi^\top\boldsymbol\Sigma\boldsymbol\Phi=\mathbf J,
  \label{eq:tw-power-preserve}
\end{equation}
即它把端口的“力--速度”功率结构精确搬运为波域的“入射--反射”功率结构。
\end{proposition}

\begin{insight}
\cref{eq:tw-power-preserve} 与哈密顿力学的正则变换同构——正则变换保持 Poisson 括号 $\{q,p\}=1$，散射变换保持功率恒等式 $f\dot x=\tfrac12(a^2-r^2)$。这并非巧合：二者都根植于能量守恒的同一数学结构（辛/双线性形式的保持）。把 $\boldsymbol\Sigma$ 看作把功率“对角化”为正、负两支的工具，散射变换就是让这两支解耦的坐标。
\end{insight}

\paragraph{从混合参数 $\mathbf H$ 到散射矩阵 $\mathbf S$：是线性分式变换，不是相似变换。}
前面 \cref{ch:teleop-twoport} 的二端口用混合参数 $\mathbf y=\mathbf H\mathbf x$，$\mathbf x=[V_m,F_e]^\top$、$\mathbf y=[F_h,-V_s]^\top$（沿用其记号）。注意 $\mathbf H$ 把\emph{混合因果}的变量相连，不能直接当作同一向量空间里的算子做相似变换。须先按“入射/反射波”重新分组：令两端口波阻抗 $b_m,b_s>0$、$\rho_m=\sqrt{2b_m}$、$\rho_s=\sqrt{2b_s}$，端口 1 流入速度 $\nu_1=V_m$、端口 2 流入速度 $\nu_2=-V_s$，则
\begin{equation}
  a_1=\frac{F_h+b_mV_m}{\rho_m},\quad r_1=\frac{F_h-b_mV_m}{\rho_m},\quad
  a_2=\frac{F_e-b_sV_s}{\rho_s},\quad r_2=\frac{F_e+b_sV_s}{\rho_s}.
  \label{eq:tw-2port-waves}
\end{equation}
把它写成 $\mathbf a=\mathbf A\mathbf x+\mathbf B\mathbf y$、$\mathbf r=\mathbf C\mathbf x+\mathbf D\mathbf y$（系数阵均为对角，元素由 \cref{eq:tw-2port-waves} 读出），代入 $\mathbf y=\mathbf H\mathbf x$ 得 $\mathbf a=(\mathbf A+\mathbf B\mathbf H)\mathbf x$、$\mathbf r=(\mathbf C+\mathbf D\mathbf H)\mathbf x$。于是在 $\mathbf A+\mathbf B\mathbf H$ 可逆时
\begin{equation}
  \boxed{\ \mathbf S=(\mathbf C+\mathbf D\mathbf H)(\mathbf A+\mathbf B\mathbf H)^{-1}\ }
  \label{eq:tw-H2S}
\end{equation}
这是 $\mathbf H\to\mathbf S$ 的 Cayley/线性分式变换。

\begin{pitfall}[把 $\mathbf H\to\mathbf S$ 当成普通相似变换]
常见错误是写 $\mathbf S=\boldsymbol\Phi\mathbf H\boldsymbol\Phi^{-1}$。这忽略了 $\mathbf H$ 的输入/输出变量是\emph{混合}的（$\mathbf x,\mathbf y$ 各取一半力、一半速度），也忽略了入射/反射波须按 \cref{eq:tw-2port-waves} 重新分组。正确关系是\cref{eq:tw-H2S} 的线性分式变换，否则会把无源性判据算错。
\end{pitfall}

\paragraph{$\mathbf S$ 与无源性的等价。}
散射矩阵 $\mathbf S$ 把入射波映到反射波，$\mathbf r=\mathbf S\mathbf a$。二端口无源当且仅当
\begin{equation}
  \|\mathbf S\|_\infty\le1,
  \label{eq:tw-Sinf}
\end{equation}
即 $\mathbf S$ 的 $H_\infty$ 范数不超过 $1$：反射波的 $L_2$ 能量不超过入射波，系统不放大能量。这与上一章用 $\mathbf H$ 表述的 Llewellyn 准则是\emph{同一个被动性条件的两种参数化}；波域的好处是判据变成简单的范数约束，且天然容纳延迟（见下节）。

\begin{pitfall}[“散射变换是一种近似”/“$b$ 须从硬件测量”]
两个新手常见误解：其一，以为“把力/速度变成波丢了信息”——其实散射变换\emph{可逆}（\cref{eq:tw-scatter-inv}），从 $(f,\dot x)$ 到 $(a,r)$ 再回来完全精确，它只是换一种“看”同一系统的方式，正如旋转坐标系不改变物理。其二，以为 $b$ 类比 $Z_0$ 就得测量——$b$ 是\emph{自由设计参数}，决定力与速度在波中的权重，任何 $b>0$ 都满足功率恒等式；如何选见 \cref{sec:tw-impedance}。
\end{pitfall}

% ════════════════════════════════════════════════════════════════
\section{波变量与 Niemeyer--Slotine 控制律}
\label{sec:tw-wavevars}

Niemeyer 与 Slotine（1991）重新包装散射变换，给出更贴合遥操作直觉的命名与一条从定义到\emph{完整控制律}的路径\cite{niemeyer1991stable}。

\paragraph{端口方向先行。}
本章端口功率以“流入通道”为正。主端取 $\nu_m=\dot x_m$；从端物理速度 $\dot x_s$ 由主端指令流向环境，对通道而言流入速度 $\nu_s=-\dot x_s$——这正是上一章 $\mathbf H$ 参数使用 $-V_s$ 的来由。

\begin{definition}[波变量]\label{def:tw-wavevars}
固定波阻抗 $b>0$。主端发射正向波 $u_m$、接收反向波 $v_m$；从端接收正向波 $u_s$、发射反向波 $v_s$：
\begin{align}
  u_m&=\frac{f_m+b\dot x_m}{\sqrt{2b}}, & v_m&=\frac{f_m-b\dot x_m}{\sqrt{2b}},\label{eq:tw-master-waves}\\
  u_s&=\frac{f_s+b\dot x_s}{\sqrt{2b}}, & v_s&=\frac{f_s-b\dot x_s}{\sqrt{2b}}.\label{eq:tw-slave-waves}
\end{align}
\end{definition}

\begin{note}
用流入通道的从端速度 $\nu_s=-\dot x_s$ 改写从端波，会更对称：$v_s=(f_s+b\nu_s)/\sqrt{2b}$、$u_s=(f_s-b\nu_s)/\sqrt{2b}$。可见 $v_s$ 才是从端“流入通道”的波、$u_s$ 是通道“送到从端”的波。\emph{勿}把从端的 $u_s$ 误读为“从从端流入通道的入射波”。又：本章小写 $v_m,v_s$ 是\emph{反向波}，与上一章大写端口速度 $V_m,V_s$ 无关，速度一律写 $\dot x$ 或 $\nu$。
\end{note}

\paragraph{通信协议：纯延迟。}
波变量的传输规则极简——正向波延迟 $T_1$ 送到从端、反向波延迟 $T_2$ 送回主端：
\begin{equation}
  u_s(t)=u_m(t-T_1),\qquad v_m(t)=v_s(t-T_2).
  \label{eq:tw-comm}
\end{equation}

\begin{center}
\small
\begin{tikzpicture}[>=Stealth, node distance=8mm, font=\small]
  \node[draw, rounded corners, fill=main!8, minimum height=11mm, minimum width=20mm] (M) {主端 $M_m$};
  \node[draw, rounded corners, fill=second!8, minimum height=11mm, minimum width=20mm, right=46mm of M] (S) {从端 $M_s$};
  \draw[->, thick] ([yshift=3mm]M.east) -- node[above, font=\scriptsize]{$u_m\!\xrightarrow{\;延迟\,T_1\;}u_s$} ([yshift=3mm]S.west);
  \draw[<-, thick] ([yshift=-3mm]M.east) -- node[below, font=\scriptsize]{$v_m\!\xleftarrow{\;延迟\,T_2\;}v_s$} ([yshift=-3mm]S.west);
  \draw[<-, thick] (M.west) -- ++(-12mm,0) node[left, font=\scriptsize]{$f_h,\dot x_m$ 人手};
  \draw[->, thick] (S.east) -- ++(12mm,0) node[right, font=\scriptsize]{$f_e,\dot x_s$ 环境};
  \node[above=1mm of M, font=\scriptsize, main!60!black] {发 $u_m$ 收 $v_m$};
  \node[above=1mm of S, font=\scriptsize, second!60!black] {收 $u_s$ 发 $v_s$};
\end{tikzpicture}
\end{center}

\paragraph{控制律：让端口波恰好等于通信波。}
关键设计是：让主/从控制器输出的力，使端口上\emph{自然形成}的波恰好匹配通信协议要传/收的波。从端收到 $u_s$ 后，由 \cref{eq:tw-slave-waves} 第一式反解出该施加的力
\begin{equation}
  \boxed{\,f_s=\sqrt{2b}\,u_s-b\,\dot x_s\,}\qquad(\text{从端控制律}).
  \label{eq:tw-slave-law}
\end{equation}
对称地，主端控制器据收到的 $v_m$ 产生
\begin{equation}
  f_{m,\rm ctrl}=\sqrt{2b}\,v_m-b\,\dot x_m,
  \label{eq:tw-master-law}
\end{equation}
以保证主端向通道送出的正向波正是 $u_m=(f_m+b\dot x_m)/\sqrt{2b}$（这里 $f_m=f_h+f_{m,\rm ctrl}$ 为主端口合力）。从端发射的反向波随之为
\begin{equation}
  v_s=\frac{f_s-b\dot x_s}{\sqrt{2b}}=\frac{(\sqrt{2b}\,u_s-b\dot x_s)-b\dot x_s}{\sqrt{2b}}=u_s-\sqrt{2b}\,\dot x_s.
  \label{eq:tw-vs-recover}
\end{equation}

\begin{insight}
\cref{eq:tw-vs-recover} 给出从端的物理图景：从端收到一个入射波 $u_s$，任务是“吸收”它——被吸收的部分转成运动 $\dot x_s$，未吸收的部分作为反射波 $v_s$ 返回主端。这与传输线末端负载完全对应：匹配负载（$Z_e=b$）反射为零、波被全吸收，透明度最高；开路（刚墙 $Z_e\to\infty$）全反射、操作者感到强力；短路（自由空间 $Z_e=0$）全反射、操作者感到的是由波阻抗 $b$ 与时延共同造成的残余惯量（$M_{\rm eff}=bT$，见 \cref{sec:tw-impedance}）。
\end{insight}

\paragraph{三段无源分解。}
Niemeyer--Slotine 证明的骨架是\emph{分解}：把系统拆成主端、通道、从端三个子系统，各证无源，再用互联定理合拢。

\begin{theorem}[波变量遥操作的无源性\,{\rm（Niemeyer--Slotine 1991）}]\label{thm:tw-passive}
若主/从端分别采用控制律 \cref{eq:tw-master-law,eq:tw-slave-law}、通信通道为纯延迟 \cref{eq:tw-comm}，则由主端、通道、从端反馈互联而成的全系统对外部（操作者与环境）无源：存在总储能 $V_{\rm tot}=V_m+V_s+E_{\rm comm}\ge0$ 使
\begin{equation}
  V_{\rm tot}(t)\le\int_0^t\big(f_h\dot x_m-f_e\dot x_s\big)\,\mathrm d\tau.
  \label{eq:tw-total-passive}
\end{equation}
\end{theorem}
\begin{proof}[证明结构]
\textbf{主端}：储能取动能 $V_m=\tfrac12M_m\dot x_m^2$，动力学 $M_m\ddot x_m=f_h+f_{m,\rm ctrl}$。由功率恒等式 $f_m\dot x_m=\tfrac12(u_m^2-v_m^2)$ 与 \cref{eq:tw-master-law}，整理得 $\dot V_m\le f_h\dot x_m-\tfrac12(u_m^2-v_m^2)$，即主端把人手输入功率减去“送入通道的净波功率”后非正增长，故无源。\textbf{从端}：储能 $V_s=\tfrac12M_s\dot x_s^2$，动力学 $M_s\ddot x_s=f_s-f_e$。注意从端流入通道速度为 $\nu_s=-\dot x_s$，通道送给从端的功率为 $\tfrac12(u_s^2-v_s^2)$，代入 \cref{eq:tw-slave-law} 得 $\dot V_s\le\tfrac12(u_s^2-v_s^2)-f_e\dot x_s$，从端无源。\textbf{通道}：纯延迟无源，$E_{\rm comm}\ge0$，详见 \cref{sec:tw-delayproof}。\textbf{互联}：三个被动子系统的功率守恒反馈互联仍被动，这正是 \cref{thm:os-interconnect} 的结论，叠加各储能即得 \cref{eq:tw-total-passive}。
\end{proof}

\begin{note}
\textbf{与本书操作空间章的接口.}\;\cref{thm:tw-passive} 的“互联”一步不重证，直接调用 \cref{ch:operational-space} 的被动系统互联定理 \cref{thm:os-interconnect}；端口功率的定义与符号约定承 \cref{eq:os-port-power}。若通道因时变时延/丢包而变有源（\cref{sec:tw-delayproof} 末与 \cref{ch:teleop-tdpa}），可再叠加一个能量罐（\cref{thm:os-energytank}）把净输出功率约束回非负——这就是下一章 TDPA 的能量层骨架。
\end{note}

% ════════════════════════════════════════════════════════════════
\section{任意常数时延下的无源性证明}
\label{sec:tw-delayproof}

这是本章\emph{皇冠上的明珠}：通道无源性的最终条件里\emph{不含}时延 $T$。

\begin{theorem}[通道时延无源性]\label{thm:tw-channel}
在任意常数时延 $T_1,T_2\ge0$ 下，波变量通信通道 \cref{eq:tw-comm} 无源，即两端口净流入能量恒非负：
\begin{equation}
  E_{\rm comm}(t)=\int_0^t\big[P_{m\to c}(\tau)+P_{s\to c}(\tau)\big]\,\mathrm d\tau\ge0,\qquad\forall t\ge0.
  \label{eq:tw-Ecomm-def}
\end{equation}
\end{theorem}

\begin{derivation}[四步证明]
\textbf{Step 1（端口功率）}\;以“流入通道”为正，用波变量展开两端口功率。主端送入通道：$P_{m\to c}=\tfrac12(u_m^2-v_m^2)$。从端用流入速度 $\nu_s=-\dot x_s$，故从端送入通道：$P_{s\to c}=\tfrac12(v_s^2-u_s^2)$。于是
\begin{equation}
  E_{\rm comm}(t)=\tfrac12\int_0^t\big[u_m^2-v_m^2+v_s^2-u_s^2\big]\,\mathrm d\tau.
  \label{eq:tw-step1}
\end{equation}
\textbf{Step 2（代入通信关系）}\;由 $u_s(\tau)=u_m(\tau-T_1)$，换元 $\sigma=\tau-T_1$（常数时延故 $\mathrm d\sigma=\mathrm d\tau$），并用因果性（$\sigma<0$ 时系统未启动，$u_m=0$）：
\begin{equation}
  \int_0^t u_s^2(\tau)\,\mathrm d\tau=\int_{-T_1}^{t-T_1}u_m^2(\sigma)\,\mathrm d\sigma=\int_0^{t-T_1}u_m^2(\sigma)\,\mathrm d\sigma.
  \label{eq:tw-step2a}
\end{equation}
同理由 $v_m(\tau)=v_s(\tau-T_2)$：$\int_0^t v_m^2\,\mathrm d\tau=\int_0^{t-T_2}v_s^2\,\mathrm d\sigma$.
\textbf{Step 3（合并）}\;代回 \cref{eq:tw-step1}：
\begin{equation}
  E_{\rm comm}(t)=\tfrac12\!\Big[\!\int_0^t\! u_m^2-\!\int_0^{t-T_1}\! u_m^2\Big]+\tfrac12\!\Big[\!\int_0^t\! v_s^2-\!\int_0^{t-T_2}\! v_s^2\Big].
  \label{eq:tw-step3}
\end{equation}
\textbf{Step 4（非负）}\;两个差恰是尾段积分：
\begin{equation}
  \boxed{\,E_{\rm comm}(t)=\tfrac12\!\int_{t-T_1}^{t}\! u_m^2(\tau)\,\mathrm d\tau+\tfrac12\!\int_{t-T_2}^{t}\! v_s^2(\tau)\,\mathrm d\tau\ \ge\ 0\,}
  \label{eq:tw-Ecomm-final}
\end{equation}
被积函数是平方项、恒非负，故积分恒非负。$\hfill\square$
\end{derivation}

\paragraph{为何 $T$ 不进入非负条件——深析。}
被积函数 $u_m^2,v_s^2\ge0$，无论积分区间 $[t-T,t]$ 多长，积分恒非负。\textbf{$T$ 只决定通道储能的\emph{大小}（区间宽度），不改变其\emph{符号}（非负性）}。物理上，\cref{eq:tw-Ecomm-final} 正是“在途波能量”——传输线里正在传播、尚未抵达的能量：

\begin{table}[H]\centering\small
\caption{通道储能两项的物理含义（在途波能量）}
\label{tab:tw-pipe-energy}
\begin{tabular}{ll}
\toprule
项 & 物理含义 \\
\midrule
$\tfrac12\int_{t-T_1}^{t}u_m^2\,\mathrm d\tau$ & 主端已发、从端尚未收到的正向波能量 \\
$\tfrac12\int_{t-T_2}^{t}v_s^2\,\mathrm d\tau$ & 从端已发、主端尚未收到的反向波能量 \\
\bottomrule
\end{tabular}
\end{table}

\begin{insight}
逆向工程看这个证明就懂它为何如此简洁：我们\emph{想要}通道储能能写成“非负被积函数之积分”。约束是变换须保持功率恒等式 $f\dot x=g(w_1,w_2)$，且 $g$ 在延迟代入后仍非负。\emph{解}就是选 $g=(u^2-v^2)/2$——功率是两个\emph{平方}之差，而平方项既非负、又在延迟下保持非负。若改选 $g=w_1w_2$，通道储能会冒出 $\int w_1(t)w_2(t-T)\,\mathrm d\tau$ 这类\emph{交叉项}，两个不同时刻信号之积可正可负，积分无法保号。散射变换的精妙正在于：把功率从“交叉项（力$\times$速度）”转成“自身平方之差（波平方差）”，消去全部交叉项。
\end{insight}

\begin{pitfall}[“波变量消除了时延的影响”]
波变量保证\emph{稳定性}，但时延仍\emph{损害透明度}：延迟到达的波解码出的力/速度是过时的，传递阻抗含 $\mathrm e^{-2\mathrm j\omega T}$ 项，高频透明度严重下降（\cref{sec:tw-impedance}）。正确理解是：波变量让你在“不稳定”与“稳定但有黏滞/软墙”之间选了后者；透明度--时延的权衡依旧存在，只是不再威胁稳定。
\end{pitfall}

\paragraph{与 Llewellyn 的对照，及证明的通用性。}
上一章 Llewellyn 准则只适用线性时不变（LTI）终端；而 \cref{thm:tw-channel} 的证明只用到功率恒等式与通信关系，操作者与环境可以是\emph{任意非线性被动}系统——这比 Llewellyn 更通用。

\begin{table}[H]\centering\small
\caption{Llewellyn 准则 vs.\ 波变量无源性}
\label{tab:tw-vs-llewellyn}
\begin{tabular}{lll}
\toprule
维度 & Llewellyn（\cref{ch:teleop-stability}） & 波变量（本章） \\
\midrule
适用范围 & 线性时不变 & 任意非线性被动终端 \\
对延迟 & 延迟使 $\eta\downarrow$，须加阻尼 & 延迟不影响无源，无需额外阻尼 \\
透明度代价 & 阻尼直接降透明度 & 波阻抗 $b$ 引入残余惯量 $bT$ \\
工程角色 & 分析工具（判稳） & 设计工具（造稳定通道） \\
\bottomrule
\end{tabular}
\end{table}

\paragraph{常数时延之外的隐忧（转交下一章）。}
\cref{thm:tw-channel} 锁定\emph{常数}时延。实际网络时延 $T(t)$ 时变：对正向波储能求导会冒出关键因子 $1-\dot T_1$；当 $\dot T_1>1$（延迟增速超过时间流逝），$-\tfrac12 u_m^2(t-T_1)(1-\dot T_1)>0$ 是额外正项——通道“压缩”信号包、功率密度被放大，能量\emph{创生}，无源性破坏。修复须在能量层显式约束输出（增益调度、能量罐 \cref{thm:os-energytank}、或时域无源 TDPA），这是 \cref{ch:teleop-tdpa} 的主题。本章到此守住“常数时延天然无源”这一最干净的结论。

% ════════════════════════════════════════════════════════════════
\section{波阻抗 \texorpdfstring{$b$}{b}：手感的旋钮}
\label{sec:tw-impedance}

$b$ 是\emph{自由设计参数}，决定力与速度在波中的相对权重（$b$ 小则 $u\approx f/\sqrt{2b}$ 偏载力，$b$ 大则 $u\approx\sqrt{b/2}\,\dot x$ 偏载速度）。它把抽象的无源性，落成“墙有多硬、自由空间有多黏”这类可感量。

\paragraph{传递阻抗的频域形。}
设环境阻抗 $Z_e(\mathrm j\omega)$、单程时延 $T$。从端反射系数、往返延迟、回主端的等效阻抗依次为
\begin{equation}
  \Gamma_s=\frac{Z_e-b}{Z_e+b},\quad \Gamma_m=\Gamma_s\,\mathrm e^{-2\mathrm j\omega T},\quad
  Z_{\rm to}(\mathrm j\omega)=b\,\frac{1+\Gamma_m}{1-\Gamma_m}=b\,\frac{1+\Gamma_s\mathrm e^{-2\mathrm j\omega T}}{1-\Gamma_s\mathrm e^{-2\mathrm j\omega T}}.
  \label{eq:tw-Zto}
\end{equation}
三种极端环境给出本章最实用的三个公式：

\begin{derivation}[刚墙：渲染刚度 $K_{\rm eff}=b/T$]
刚墙 $Z_e\to\infty$ 故 $\Gamma_s\to+1$（同相全反射，亦可由边界条件 $\dot x_s=0$ 代入 \cref{eq:tw-vs-recover} 得 $v_s=u_s$ 看出）。代入 \cref{eq:tw-Zto} 并用 Euler 公式：
\begin{equation}
  Z_{\rm to}=b\,\frac{1+\mathrm e^{-2\mathrm j\omega T}}{1-\mathrm e^{-2\mathrm j\omega T}}
          =\frac{b}{\mathrm j\tan(\omega T)}\ \xrightarrow{\ \omega T\ll1\ }\ \frac{b}{\mathrm j\omega T}=\frac{K_{\rm eff}}{\mathrm j\omega},\quad K_{\rm eff}=\frac{b}{T}.
  \label{eq:tw-Keff}
\end{equation}
\end{derivation}

低频下波变量把刚墙渲染为有限刚度 $K_{\rm eff}=b/T$：时延越大或 $b$ 越小，墙越“软”。典型地 $b=2\,$Ns/m、$T=100\,$ms 给 $K_{\rm eff}=20\,$N/m，而真墙刚度可达 $10^4\,$N/m——渲染失真约 $500$ 倍，这正是“高时延下墙像橡胶”的定量根源。

\begin{derivation}[自由空间：残余惯量 $M_{\rm eff}=bT$]
自由空间 $Z_e=0$ 故 $\Gamma_s\to-1$，代入 \cref{eq:tw-Zto}：
\begin{equation}
  Z_{\rm to}=b\,\frac{1-\mathrm e^{-2\mathrm j\omega T}}{1+\mathrm e^{-2\mathrm j\omega T}}=b\,\mathrm j\tan(\omega T)\ \xrightarrow{\ \omega T\ll1\ }\ \mathrm j\omega\,bT=\mathrm j\omega\,M_{\rm eff},\quad M_{\rm eff}=bT.
  \label{eq:tw-Beff}
\end{equation}
\end{derivation}

自由空间里操作者感到的不是阻尼，而是一团额外\emph{惯量} $M_{\rm eff}=bT$——传递阻抗低频展开为 $Z_{\rm to}\approx\mathrm j\omega\,bT$，与质量 $f=M\ddot x$ 的阻抗 $\mathrm j\omega M$ 同型（阻尼 $f=B\dot x$ 的阻抗是\emph{频率无关}的 $B$，二者形式不同）。这团随时延线性增大的“附加质量”，正是波变量在自由空间里那种“推着重物、起停发沉”的迟滞手感的来源，文献中亦称之为延迟诱导的 added inertia/apparent mass。第三种情况是\emph{完美匹配} $Z_e=b$，则 $\Gamma_s=0$、$Z_{\rm to}=b$ 在\emph{所有}频率成立——此时操作者感到的是\emph{纯阻尼} $b$（反射最小、透明度最高，但所有频率都遮蔽为 $b$）。

\begin{note}
\textbf{临界频率与共振.}\;在 $\omega=\pi/(2T)$ 处 $\tan(\omega T)\to\infty$：刚墙 $Z_{\rm to}\to0$（反共振），自由空间 $Z_{\rm to}\to\infty$（共振），透明度在该频率彻底丧失。故波变量的有效带宽约 $\omega<1/(2T)$。
\end{note}

\paragraph{阻抗匹配、振铃与综合权衡。}
$K_{\rm eff}=b/T$（要墙硬，须大 $b$）与 $M_{\rm eff}=bT$（要自由空间轻，须小 $b$）方向相反——这就是 $b$ 选取的核心张力。当 $b\ll Z_e$ 时失配比 $\Gamma\approx1$、近全反射，反射波在主从间来回弹跳、逐次被部分吸收，形成\emph{振铃}，衰减时间常数 $\tau_{\rm ring}\approx T/\ln(1/|\Gamma|)\approx TZ_e/(2b)$，$b$ 太小时可长达数秒。综合最优是\emph{阻抗匹配} $b=\sqrt{K_eM_m}$：

\begin{insight}
$b=\sqrt{K_eM_m}$ 有两重等价解读。\emph{传输线视角}：特征阻抗等于负载阻抗时反射为零、传输效率最高（消除驻波）。\emph{最优控制视角}：它是 $\min_b\int_0^\infty|\Gamma(\omega)|^2\,\mathrm d\omega$ 的解，即“所有频率反射能量之和”最小。两条路通向同一点，因为它们刻画的是同一个物理——能量不被反射、不在管线里空转。
\end{insight}

\begin{table}[H]\centering\small
\caption{波阻抗 $b$ 的定量权衡（以 $\sqrt{K_eM_m}$ 为单位）}
\label{tab:tw-b-tradeoff}
\begin{tabular}{lllll}
\toprule
$b/\sqrt{K_eM_m}$ & 自由空间附加惯量 & 硬墙振铃 & 阻抗匹配度 & 评价 \\
\midrule
$0.1$ & 极低（好） & 严重（差） & $10\%$ & 不可用 \\
$0.3$ & 低（好） & 明显 & $30\%$ & 自由空间任务可用 \\
$0.5$ & 中等 & 轻微 & $50\%$ & 可接受 \\
$\mathbf{1.0}$ & 中等 & \textbf{最小} & $\mathbf{100\%}$ & \textbf{最优匹配} \\
$2.0$ & 高（差） & 无 & $50\%$（过阻尼） & 保守 \\
$5.0$ & 严重（差） & 无 & $20\%$（严重过阻尼） & 不可用 \\
\bottomrule
\end{tabular}
\end{table}

\paragraph{自适应 $b$。}
环境刚度未知或时变时，可在线估计 $\hat Z_e=|f_e/\dot x_s|$ 并令 $b\!\leftarrow\!\sqrt{\hat Z_e M_m}$，经低通滤波后限幅。\textbf{关键}：$b$ 的变化\emph{必须平滑}——突变会打破波变量连续性、注入能量脉冲（等效于在波域引入一个非物理阶跃）。典型经验范围：平动 $b\in[0.5,5]\,$Ns/m（接近人手操作阻抗量级）、转动 $b\in[0.05,0.5]\,$Nms/rad（接近手腕阻抗量级）。

\begin{pitfall}[“$b$ 越大越稳所以越好”]
大 $b$ 确实压制振铃，但代价是严重的附加惯量（$M_{\rm eff}=bT$）——自由空间起停发沉、“像推着重物”，且力信息被速度信息淹没。$b$ 的最优值取决于环境刚度，不是“越大越好”。把 $b$ 当作\emph{透明度--鲁棒性的调节旋钮}，而非“安全系数”。
\end{pitfall}

% ════════════════════════════════════════════════════════════════
\section{位置漂移与波积分（指向下一章）}
\label{sec:tw-drift}

\paragraph{漂移从何而来。}
波变量传输的是\emph{速度}信息（编码在波中），位置须在接收端\emph{积分}恢复。从端速度 $\dot x_s=(u_s-v_s)/\sqrt{2b}$ 积分得位置，而积分器 $1/s$ 在零频有极点、对直流偏置无界放大：任何微小直流偏置 $\delta$（量化、零偏、舍入）都使位置误差线性增长
\begin{equation}
  \Delta x(t)=\frac{\delta}{\sqrt{2b}}\,t.
  \label{eq:tw-drift}
\end{equation}
即便 $\delta=10^{-5}\sqrt{\text W}$、$b=2$，一小时后 $\Delta x\approx18\,$mm——已可感知。丢包会进一步使速度采样缺失、漂移按丢包率加剧。

\begin{table}[H]\centering\small
\caption{位置漂移的来源}
\label{tab:tw-drift-sources}
\begin{tabular}{lll}
\toprule
来源 & 机制 & 量级 \\
\midrule
初始化 & $x_m(0)\neq x_s(0)$ & 固定偏移 \\
数值积分 & 累积舍入误差 & 随时间线性增长 \\
丢包 & 速度采样丢失 & 随丢包率增加 \\
$b$ 切换 & 自适应 $b$ 瞬间不连续 & 脉冲偏移 \\
\bottomrule
\end{tabular}
\end{table}

\paragraph{波积分（wave integral）：思路与无源代价。}
Niemeyer 的修复是同时传输波的\emph{积分} $U_m(t)=\int_0^t u_m\,\mathrm d\tau$。由通信关系 $U_s(t)=U_m(t-T_1)$，从端可重建“位置指令”
\begin{equation}
  x_s^{\rm cmd}(t)=\frac{U_s(t)-V_s(t)}{\sqrt{2b}}=\frac{U_m(t-T_1)-V_s(t)}{\sqrt{2b}},
  \label{eq:tw-xcmd}
\end{equation}
再用一根\emph{弱弹簧} $f_{\rm corr}=K_{\rm drift}(x_s^{\rm cmd}-x_s)$ 把漂移拉回。弹簧是被动元件，储能 $\tfrac12K_{\rm drift}(x_s^{\rm cmd}-x_s)^2\ge0$、只存不创，故不破坏 \cref{thm:tw-passive} 的无源性。$K_{\rm drift}$ 须\emph{弱}（典型 $10$--$50\,$N/m，远小于环境刚度）以不损透明度，但受 Z-width 约束 $K_{\rm drift}<2b/T_s$（采样周期 $T_s$，承上一章离散无源边界）。完整工程实现（含 double 精度积分、自由空间“回飘”抑制）属下游内容，留 \cref{ch:teleop-tdpa}。

\begin{insight}
波积分的漂移修复与 SLAM 里的回环修正同源：纯里程计积分因累积误差漂移、回环检测提供约束把它拉回；波积分里传输的 $U_m$ 提供“参考位置约束”、弱弹簧把积分漂移拉回。二者都是“外部参考 $+$ 弱反馈”修复积分器固有的不稳定性——把纯积分 $1/s$ 改成一阶低通 $1/(s+K_{\rm drift}/b)$（弱弹簧 $K_{\rm drift}$ 与端口在自由空间呈现的阻尼 $b$ 一起把极点从原点拉到 $-K_{\rm drift}/b$）。
\end{insight}

% ════════════════════════════════════════════════════════════════
\section{本章小结}
\label{sec:tw-summary}

\paragraph{主线回顾。}
本章把遥操作的时延稳定性从\emph{控制问题}重述为\emph{信号表示问题}：延迟保能、放大才产能（\cref{sec:tw-motivation}）$\to$ 无损传输线的功率分解 $P=(V^+)^2/Z_0-(V^-)^2/Z_0$ 提供原型（\cref{sec:tw-transmissionline}）$\to$ 散射变换把端口功率对角化为波平方差 $f\dot x=\tfrac12(a^2-r^2)$、并以 $\boldsymbol\Phi^\top\boldsymbol\Sigma\boldsymbol\Phi=\mathbf J$ 与 $\|\mathbf S\|_\infty\le1$ 把无源性化为范数约束（\cref{sec:tw-scattering}）$\to$ 波变量给出对称的主/从控制律 $f_s=\sqrt{2b}\,u_s-b\dot x_s$，三段分解复用互联定理证全系统无源（\cref{sec:tw-wavevars}）$\to$ 四步证明锁定“任意常数时延天然无源、$T$ 不入非负条件”这一皇冠结论（\cref{sec:tw-delayproof}）$\to$ 波阻抗 $b$ 把抽象无源落成 $K_{\rm eff}=b/T$、$M_{\rm eff}=bT$、匹配 $b=\sqrt{K_eM_m}$ 等可感量（\cref{sec:tw-impedance}）$\to$ 残留的位置漂移由波积分修复、转交时变时延与 TDPA（\cref{sec:tw-drift}）。

\paragraph{与他章的关系。}
向上承 \cref{ch:teleop-twoport}（二端口/混合参数）与 \cref{ch:teleop-stability}（Llewellyn/时延使 $\eta<0$）：本章给出绕过透明度--稳定性矛盾的结构性解法。横向复用 \cref{ch:operational-space} 的端口功率 \cref{eq:os-port-power}、耗散不等式 \cref{eq:os-passive-integral}、互联定理 \cref{thm:os-interconnect}、能量罐 \cref{thm:os-energytank}，不重证系统级无源性理论。向下，时变时延/丢包的能量层修复（增益调度、能量罐、时域无源）与波积分工程化在 \cref{ch:teleop-tdpa} 展开；遥操作的运动映射（主从工作空间标度 $x_s=\lambda_p x_m$、奇异规避）见 \cref{ch:teleop-mapping}；接触安全与碰撞处理见 \cref{ch:arm-collision-safety}。

\begin{table}[H]\centering\small
\caption{本章速查：散射变换、波变量、时延无源与波阻抗}
\label{tab:tw-cheatsheet}
\begin{tabular}{p{0.34\textwidth} p{0.58\textwidth}}
\toprule
条目 & 关键式 / 判据 \\
\midrule
散射波定义 & $a=(f+b\dot x)/\sqrt{2b}$，$r=(f-b\dot x)/\sqrt{2b}$ \\
逆变换 & $f=\sqrt{b/2}\,(a+r)$，$\dot x=(a-r)/\sqrt{2b}$ \\
功率恒等式 & $f\dot x=\tfrac12(a^2-r^2)$ \\
散射矩阵 & $\boldsymbol\Phi=\tfrac1{\sqrt{2b}}\!\big[\begin{smallmatrix}1&b\\1&-b\end{smallmatrix}\big]$，$\det=-1$ \\
功率保持 & $\boldsymbol\Phi^\top\boldsymbol\Sigma\boldsymbol\Phi=\mathbf J$ \\
$\mathbf H\!\to\!\mathbf S$ & $\mathbf S=(\mathbf C+\mathbf D\mathbf H)(\mathbf A+\mathbf B\mathbf H)^{-1}$（线性分式，非相似变换） \\
无源等价 & $\|\mathbf S\|_\infty\le1$ \\
通信协议 & $u_s(t)=u_m(t-T_1)$，$v_m(t)=v_s(t-T_2)$ \\
从端控制律 & $f_s=\sqrt{2b}\,u_s-b\dot x_s$，$v_s=u_s-\sqrt{2b}\,\dot x_s$ \\
通道储能 & $E_{\rm comm}=\tfrac12\!\int_{t-T_1}^t\! u_m^2+\tfrac12\!\int_{t-T_2}^t\! v_s^2\ge0$（$T$ 不入非负条件） \\
刚墙渲染刚度 & $K_{\rm eff}=b/T$ \\
自由空间残余惯量 & $M_{\rm eff}=bT$（$Z_{\rm to}=\mathrm j\omega bT$，质量型） \\
阻抗匹配 & $b=\sqrt{K_eM_m}$，$\Gamma=(Z_e-b)/(Z_e+b)$ \\
有效带宽 & $\omega<1/(2T)$；$\omega=\pi/(2T)$ 处透明度尽失 \\
时变时延破坏 & 关键因子 $1-\dot T_1$；$\dot T_1>1$ 时能量创生 \\
位置漂移修复 & 波积分 $x_s^{\rm cmd}=(U_m(t-T_1)-V_s)/\sqrt{2b}$ $+$ 弱弹簧 $K_{\rm drift}$ \\
\bottomrule
\end{tabular}
\end{table}

\subsection*{常见误解汇总}
\begin{itemize}
  \item \textbf{“波变量让延迟消失”}：延迟仍在、透明度仍降；波变量只让延迟从不安全变安全（\cref{sec:tw-motivation}）。
  \item \textbf{“散射变换是近似/丢信息”}：它是可逆线性变换，往返完全精确（\cref{eq:tw-scatter-inv}）。
  \item \textbf{“$\mathbf H\to\mathbf S$ 是相似变换”}：是线性分式变换 \cref{eq:tw-H2S}，因 $\mathbf H$ 变量混合因果。
  \item \textbf{“$b$ 须从硬件测量”}：$b$ 是自由设计参数，任何 $b>0$ 都满足功率恒等式（\cref{sec:tw-impedance}）。
  \item \textbf{“常数时延无源 $\Rightarrow$ 时变时延也无源”}：$\dot T>1$ 时能量创生，须能量层修复（\cref{ch:teleop-tdpa}）。
\end{itemize}

\subsection*{符号表}
\begin{center}\small
\begin{tabular}{ll}
\toprule
符号 & 含义 \\
\midrule
$f,\ \dot x$ & 端口力、速度（端口功率 $P=f\dot x$，承 \cref{eq:os-port-power}） \\
$b$ & 波阻抗（自由设计参数，$>0$） \\
$a,\ r$ & 单端口入射波、反射波（\cref{eq:tw-scatter-def}） \\
$u_m,v_m,u_s,v_s$ & 主/从端正向波、反向波（\cref{def:tw-wavevars}） \\
$\nu_m,\nu_s$ & 流入通道的端口速度，$\nu_m=\dot x_m$、$\nu_s=-\dot x_s$ \\
$T_1,T_2,\ T$ & 正/反向单程时延、往返时延 \\
$\boldsymbol\Phi,\ \mathbf S$ & 单端口散射矩阵、二端口散射矩阵 \\
$\mathbf J,\ \boldsymbol\Sigma$ & 端口功率型、波功率型矩阵（\cref{eq:tw-JSigma}） \\
$Z_e,\ Z_{\rm to},\ \Gamma$ & 环境阻抗、传递阻抗、反射系数 \\
$K_{\rm eff},\ M_{\rm eff}$ & 刚墙渲染刚度 $b/T$、自由空间残余惯量 $bT$（质量型，$Z_{\rm to}=\mathrm j\omega bT$） \\
$E_{\rm comm}$ & 通道（在途波）储能，恒非负 \\
$K_{\rm drift}$ & 波积分弱弹簧增益 \\
$M_m,\ M_s$ & 主端、从端等效惯量 \\
\bottomrule
\end{tabular}
\end{center}

\subsection*{延伸阅读}
散射变换原始文献为 Anderson--Spong（1989）\cite{anderson1989bilateral}；波变量与“稳定自适应遥操作”为 Niemeyer--Slotine（1991）\cite{niemeyer1991stable}，位置漂移与波积分（wave integral）的系统处理见 Niemeyer 的 MIT 博士论文（1996，\emph{Using Wave Variables in Time Delayed Force Reflecting Teleoperation}）及其期刊版 Niemeyer--Slotine（2004，\emph{Telemanipulation with Time Delays}, IJRR 23(9):873--890，明确给出“既不漂移又无源”的传输方案）；无源遥操作的系统综述见 Nu\~no--Basa\~nez--Ortega（2011）\cite{nuno2011passivity}；时域无源控制（TDPA，下一章主线）为 Hannaford--Ryu（2002）\cite{hannaford2002time}；耦合稳定与无源性的奠基为 Colgate--Hogan（1988）\cite{colgate1988robust}；几何/端口视角的无源框架见 Stramigioli（2001）\cite{stramigioli2001modeling}。时变时延的（时变增益）无源化修复见 Lozano--Chopra--Spong（Mechatronics'02, Enschede, 2002）；能量罐/透明度的双层架构见 Franken 等（IEEE T-RO 27(4):741--756, 2011）。
